\chapter{Conclusion}

The repository now contains enough material to support a genuine monograph-level thesis.
The thesis is not merely that sheaf theory is a pretty metaphor for program analysis, nor
merely that hybrid dependent types are a convenient wrapper around solver checks and LLM
judgments. The deeper thesis is that local-to-global semantics can be a universal
organizing principle across a surprising amount of software engineering work.

The plain type system already provides a strong foundational language of exact semantic
objects. The hybrid system already provides a layered semantics of trust-bearing evidence.
The geometry modules already provide a language of gluing and obstruction. The equivalence
modules already show how local isomorphism can become a global verdict. The anti-
hallucination modules already show how artifact quality can be recast as inhabitation.
Type expansion and the agent already show that planning and generation can be read as typed
elaboration. Zero-change checking and mixed mode already show that the architecture can be
adopted both from legacy Python and from proof-oriented new code.

What the revised manuscript has tried to add is not merely more detail but more inevitability.
It has tried to show that these lines of work do not just coexist; they compel one another.
Once one decides that software facts should be represented locally, that evidence should be
typed, that differences in evidence quality should remain explicit, and that non-gluing
should be treated as structure rather than noise, then a remarkable amount follows. One is
led toward refinement lattices, toward hybrid witnesses, toward trust-sensitive semantics,
toward cohomological diagnosis, toward theory packs, toward zero-change extraction, and
toward synthesis as elaboration. The unity of the project lies not in a branding exercise,
but in this chain of implication.

There is still real work ahead. The repository's different lines of development need to be
made even more consistent. Trust vocabularies need consolidation. The bridge from intent to
formalization remains heuristic and should be described as such. Covers, site families, and
residuals should be made still more explicit in user-facing outputs. But those are exactly
the kinds of tasks one expects in a large research prototype converging on a unifying
architecture.

This is also why the phrase ``absolutely everything'' should be read with discipline rather
than skepticism. It does not mean that every domain has already been mastered. It means that
the project has identified an extensible semantic shape into which new domains can be fit.
Whenever one can identify local sites, typed local evidence, transport maps, and a sensible
notion of gluing, one has a candidate entrance into the same universe. The right way to
extend the system is therefore not to bolt on one more isolated feature, but to ask what the
local sections are, where their overlaps live, what trust structure they require, and how
their obstructions should be repaired.

If the phrase ``cohomological typing for absolutely everything'' sounded excessive at the
start, the point of this monograph is to make it sound less like excess and more like a
research program. The program is this: whenever a subsystem of software work can be
understood as local evidence that should glue globally, \DepPy{} asks whether typing,
trust, and cohomology can be used to express that subsystem in one semantic language. Much
of the repository is already an existence proof that the answer is often yes.

The deepest hope behind the project is therefore not just that particular theorems will be
proved or that particular agents will generate better code. It is that software engineering
can acquire a more stable conceptual center. In this monograph, that center is the claim
that types, evidence, proofs, semantics, and generation are not unrelated vocabularies but
progressive refinements of one local-to-global discipline. If the reader leaves with that
impression, the document will have done what a foundational document should do: make later
work feel less optional and more like the next chapter of the same idea.
